%
% Trabalho de Conclusão de Curso (TCC)
% Campus Belo Jardim
% Instituto Federal de Pernambuco (IFPE)
% Belo Jardim Pernambuco Brasil
% 18 de Outubro de 2014
%
% Autor: 

%

\documentclass[oneside, 12pt]{book}

\usepackage[T1]{fontenc}
\usepackage{quoting}
\usepackage[brazil]{babel}
\usepackage[utf8]{inputenc}
\usepackage{amsfonts}
\usepackage{amssymb}
\usepackage{graphicx}
\usepackage{units}
\usepackage{makeidx}  
\usepackage{float}
%\usepackage[nottoc]{tocbibind}  
\usepackage{placeins}
\usepackage{amsmath}
\usepackage{tabularx,colortbl}
\usepackage{multirow}
\usepackage{color}
\usepackage{xcolor}
\usepackage{url}
\usepackage{lscape}
\usepackage{subcaption}
\usepackage{caption}
\usepackage{setspace} 
\usepackage[a4paper,top=3cm,bottom=2.0cm,left=3.0cm,right=2cm]{geometry} 
\usepackage[titletoc]{appendix}
\usepackage{longtable}
\usepackage{titling}
\usepackage{pdfpages}
\usepackage{titlesec}

\usepackage[ruled,vlined,portuguese]{algorithm2e}
\SetKwFor{ParaTodo}{para todo}{fa\c{c}a}{fim para todo}
\SetKwBlock{Inicio}{início}{fim}
%\usepackage{algorithmicx}
\usepackage{algpseudocode}

% ---------------------------------------------------------------------------- %
% Cabeçalhos similares ao TAOCP de Donald E. Knuth
\usepackage{fancyhdr}
\pagestyle{fancy}
\fancyhf{}
\titleformat{\chapter}[display]
  {\normalfont\huge\bfseries}{\thechapter}{20pt}{\Huge}
\renewcommand{\chaptermark}[1]{\markboth{\MakeUppercase{#1}}{}}
\renewcommand{\sectionmark}[1]{\markright{\MakeUppercase{#1}}{}}

\renewcommand{\headrulewidth}{0pt}
\setlength{\headheight}{15.1pt}

\frenchspacing                          % arruma o espaço: id est (i.e.) e exempli gratia (e.g.) 
\urlstyle{same}                         % URL com o mesmo estilo do texto e não mono-spaced
\makeindex                              % para o índice remissivo
\raggedbottom                           % para não permitir espaços extra no texto
\cleardoublepage
\normalsize

\newcolumntype{C}[1]{>{\centering\arraybackslash}p{#1}}

\usepackage{csquotes}
\usepackage[style=abnt]{biblatex}
\bibliography{bibliography}
\addbibresource{bibliography.bib}

%--------------------------------------------
% Variáveis importantes
\title{Plataforma de Engenharia Self-Service: Abstraindo a Complexidade da Infraestrutura para Pesquisa e Inovação}
\author{Fábio Braz Ferreira}
\date{DD/MM/AAAA}
% Espaço reservado para figura: \figplaceholder{rótulo}{legenda}{o que inserir}{fonte}
% Para trocar pela imagem real, substitua o \fbox por \includegraphics[width=0.9\textwidth]{images/arquivo.png}
\newcommand{\figplaceholder}[4]{%
  \begin{figure}[htb]
    \centering
    \fbox{\parbox[c][5cm][c]{0.85\textwidth}{\centering\color{gray}[INSERIR IMAGEM]\\[2mm]\footnotesize #3}}
    \caption{#2}
    \label{#1}
    \par\small Fonte: #4.
  \end{figure}%
}

% Figura real: \figimg{rótulo}{legenda}{arquivo em images/ (sem extensão)}{fonte}
\newcommand{\figimg}[4]{%
  \begin{figure}[htb]
    \centering
    \includegraphics[width=0.95\textwidth]{images/#3}
    \caption{#2}
    \label{#1}
    \par\small Fonte: #4.
  \end{figure}%
}

\newcommand{\thecourse}{Bacharelado em Engenharia de Software}
\newcommand{\thesupervisor}{Elton Bezerra Torres}
\newcommand{\thelocal}{Belo Jardim, Pernambuco}
\newcommand{\thequalifiers}{
{Elton Bezerra Torres - IFPE - Campus Belo Jardim (Orientador)} \hfill  {}\\
{NOME DO CONVIDADO 1 - IFPE - Campus Belo Jardim} \hfill  {}\\
{NOME DO CONVIDADO 2 - INSTITUIÇÃO}
}

\begin{document}
	\frontmatter
    \renewcommand{\refname}{Referências}% para classes report, article, etc
    \renewcommand{\bibname}{Referências}% para classes baseadas no book
	
	\setcounter{tocdepth}{2}
	%\fancyhead[LE]{\thepage\hspace{2em}\footnotesize{\leftmark}}
	\fancyhead[RO]{{\footnotesize\rightmark}\hspace{2em}\thepage}
	
	\begin{titlepage}
	\begin{figure}[ht]
        \centering
        \includegraphics[width = 0.25\textwidth]{images/brazao_rep.jpg}
	\end{figure}
 \centering
         \textbf{INSTITUTO FEDERAL DE PERNAMBUCO}\\
        \textbf{\textit{CAMPUS} BELO JARDIM}\\
        \textbf{BACHARELADO EM ENGENHARIA DE SOFTWARE}\\
	\begin{center}

        \vspace{2in}
		{\textbf{\MakeUppercase{\theauthor}}}\\[3cm]
		
		\Large{\textbf{\textbf{\thetitle}}}\\[1.5cm]
		\vspace{1.5in}
		\large{\textbf{\thelocal}\\
			{\thedate}}
	\end{center}
\end{titlepage}



	\newpage
	\begin{titlepage}
	
\begin{center}

%==(CABEÇALHO)========================================================

{\textbf{\MakeUppercase{\theauthor}} \\ \vspace{1ex}

%=====================================================================

\vspace{1.8in}

%===(TÍTULO DO PROJETO DE DISSERTAÇÃO)================================
{\Large \textbf{\thetitle} } \\
%=====================================================================

\vspace{0.8in}
}

\end{center}

%===(IDENTIFICADOR DO DOCUMENTO)===================================
\begin{flushright}
\vspace{0.5in}
\parbox{2.75in}
{\textbf{Trabalho de conclusão de Curso (TCC)} apresentado como requisito parcial para obtenção do grau de Bacharel em Engenharia de Software. Primeiro TCC aprovado no curso de Engenharia de Software do IFPE - Campus Belo Jardim}
\end{flushright}
%===================================================================


%===(Banca de Qualificação)=========================================
\begin{flushright}
\textbf{Banca de Qualificação:} \\ \vspace{1ex}
\thequalifiers

\end{flushright}
%===================================================================



\vspace{1in}

\begin{center}

\vspace{50pt} {\thelocal, \thedate.}
%==================================================================

\end{center}


\end{titlepage}



	

	\newpage
    \chapter*{}
\begin{center}
    \centering {\textbf{\MakeUppercase{\theauthor}}}

\vspace{2.5in}

{\textbf{\MakeUppercase{\thetitle}}}

\vspace{0.5in}

\begin{center}
    {\centering \textbf{\thesupervisor} \par}
    \makebox[10cm]{\hrulefill}\\
    {\centering Professor Orientador \par}

    \vspace{0.5in}

    {\centering \textbf{Hítalo Oliveira da Silva} \par}
    \makebox[10cm]{\hrulefill}\\
    {\centering Convidado 1 \par}

    \vspace{0.5in}

    {\centering \textbf{Daniel Leite Viana} \par}
    \makebox[10cm]{\hrulefill}\\
    {\centering Convidado 2 \par}
\end{center}
\end{center}
    \newpage

	
	% ---------------------------------------------------------------------------- %
% Agradecimentos:
% Se o candidato não quer fazer agradecimentos, deve simplesmente eliminar esta página
\chapter*{Agradecimentos}

% TODO: personalizar este texto — o parágrafo abaixo é um ponto de partida.

A elaboração deste trabalho não teria sido possível sem o apoio e a colaboração de diversas pessoas e instituições, as quais gostaria de agradecer e reconhecer.\\

Ao Instituto Federal de Pernambuco – Campus Belo Jardim, pela formação recebida ao longo do curso de Engenharia de Software.\\

Ao meu orientador, professor Elton Bezerra Torres, pelo direcionamento e pelas contribuições ao longo do desenvolvimento deste trabalho.\\

À minha família e aos amigos, pelo apoio e incentivo durante essa etapa.\\


	\chapter*{Resumo}

\par Este trabalho apresenta o desenvolvimento de uma Plataforma de Engenharia Self-Service com o objetivo de abstrair a complexidade na provisão e configuração de ambientes de infraestrutura para pesquisa e desenvolvimento. Tradicionalmente, a criação de máquinas virtuais (VMs) ou contêineres, bem como a configuração de softwares e serviços específicos, demanda tempo e conhecimento técnico especializado do usuário final. A solução desenvolvida, denominada ProxLab, transforma esse processo manual em um modelo sob demanda: por meio de um portal web, o usuário (aluno ou pesquisador) especifica suas necessidades de infraestrutura, e o backend orquestra a criação de máquinas virtuais no Proxmox VE por meio de Terraform (Infraestrutura como Código) ou o provisionamento de contêineres de aplicação em um cluster Kubernetes (k3s) dedicado, com quotas por perfil de usuário e acesso interativo via terminal web. Essa abordagem de Engenharia de Plataforma capacita pesquisadores que não possuem conhecimento aprofundado em infraestrutura a provisionarem seus próprios ambientes de trabalho de forma ágil, reduzindo a dependência de equipes de operação e a ocorrência de erros manuais. Os resultados obtidos demonstram a viabilidade técnica da abordagem, evidenciando ganhos de autonomia e redução do tempo de provisionamento em relação ao processo manual.

\vspace{2em}

\textbf{Palavras-chave:} Engenharia de Plataforma, Self-Service, Infraestrutura como Código (IaC), Terraform, Kubernetes, Proxmox, Automação de Infraestrutura.

\vfill

% ---------------------------------------------------------------------------- %
	% ---------------------------------------------------------------------------- %
% Abstract
\chapter*{Abstract}

This work presents the development of a Self-Service Engineering Platform aimed at abstracting the complexity of provisioning and configuring infrastructure environments for research and development. Traditionally, creating virtual machines (VMs) or containers, as well as configuring specific software and services, demands time and specialized technical knowledge from the end user \cite{KunalBhandari}. The developed solution, named ProxLab, turns this manual process into an on-demand model: through a web portal, the user (student or researcher) specifies their infrastructure needs, and the backend orchestrates the creation of virtual machines on Proxmox VE through Terraform (Infrastructure as Code) or the provisioning of application containers on a dedicated Kubernetes (k3s) cluster, with per-profile quotas and access via web terminal or network. This Platform Engineering approach \cite{Sandobalín} empowers researchers who do not have in-depth infrastructure knowledge to provision their own work environments in an agile way, reducing dependency on operations teams and the occurrence of manual errors \cite{Gokhale}. The results obtained demonstrate the technical feasibility of the approach, evidencing gains in autonomy and reduction in provisioning time compared to the manual process \cite{Muhammad}.

\vfill \textbf{Keywords:} Platform Engineering, Self-Service, Infrastructure as Code (IaC), Terraform, Kubernetes, Proxmox, Infrastructure Automation.
\\

% ---------------------------------------------------------------------------- %\\
	\chapter*{Lista de Abreviaturas}
\begin{longtable}{ll}
        IaC	& Infraestrutura como Código (\emph{Infrastructure as Code}) \\
        VM	& Máquina Virtual (\emph{Virtual Machine}) \\
        API	& Interface de Programação de Aplicações (\emph{Application Programming Interface}) \\
        REST & \emph{Representational State Transfer} \\
        JWT & Token Web JSON (\emph{JSON Web Token}) \\
        ORM & Mapeamento Objeto-Relacional (\emph{Object-Relational Mapping}) \\
        CRUD & Criar, Ler, Atualizar, Deletar (\emph{Create, Read, Update, Delete}) \\
        RX & Experiência do Pesquisador (\emph{Researcher Experience}) \\
        PoC & Prova de Conceito (\emph{Proof of Concept}) \\
        SSH & \emph{Secure Shell} \\
        TLS & Segurança da Camada de Transporte (\emph{Transport Layer Security}) \\
        IFPE & Instituto Federal de Pernambuco \\
\end{longtable}

	\include{sections/01_pre_text/007_lista_simbolos}

    \listoffigures		
	\listoftables
    \renewcommand\contentsname{\centering Sumário} 
	\tableofcontents
	
	
	\mainmatter	
	% cabeçalho para as páginas de todos os capítulos
	\fancyhead[LO]{\thesection}
	
	%\singlespacing            % espaçamento simples
	\onehalfspacing            % espaçamento um e meio
	
	\setcounter{page}{1}
	\chapter{Introdução}

A crescente demanda por agilidade e eficiência no ciclo de desenvolvimento de software e pesquisa tem impulsionado transformações significativas na maneira como a infraestrutura de tecnologia é gerenciada e provisionada. Tradicionalmente, o processo de solicitação e configuração de ambientes computacionais, sejam máquinas virtuais (VMs) ou contêineres, bem como a instalação de dependências e ferramentas específicas, representa um gargalo considerável para desenvolvedores e pesquisadores. A complexidade inerente à gestão de infraestrutura de baixo nível, a necessidade de múltiplos conhecimentos técnicos e o tempo de espera envolvido na interação com equipes de infraestrutura ou operações frequentemente resultam em atrasos no projeto e na diminuição da produtividade \cite{KunalBhandari}.\\

Nesse cenário, a \textit{Engenharia de Plataforma} surge como uma disciplina fundamental. Ela se concentra em criar e manter ferramentas e infraestrutura que capacitam os usuários, sejam desenvolvedores, pesquisadores ou outros profissionais, a consumir recursos de tecnologia de forma autônoma e eficiente, sem a necessidade de aprofundar-se nos detalhes complexos da infraestrutura subjacente \cite{cncfPlatforms:online,Anjum2026}. O objetivo primário é otimizar a Experiência do Pesquisador (\textit{Researcher Experience} — RX) — adaptação do conceito de \textit{Developer Experience} \cite{Fagerholm2012,Greiler2023} —, permitindo que os pesquisadores se concentrem em suas atividades principais, e não no provisionamento e na manutenção de seus ambientes.\\

Diante desse contexto, este trabalho apresenta o desenvolvimento do \textbf{ProxLab}, uma \textbf{Plataforma de Engenharia Self-Service} voltada para alunos e pesquisadores de um ambiente acadêmico de virtualização \textit{on-premises}. A plataforma oferece um portal web onde os usuários definem suas necessidades de infraestrutura, escolhendo entre máquinas virtuais e contêineres de aplicação, dentro de limites de recursos definidos por perfil de usuário (aluno ou pesquisador). Diferentemente da abordagem inicialmente concebida — que previa a automação de configuração por meio do Ansible —, a solução final adota uma arquitetura híbrida: a orquestração de máquinas virtuais no Proxmox VE é feita via \textit{Terraform} (Infraestrutura como Código), enquanto o provisionamento de contêineres de aplicação é realizado em um cluster Kubernetes (\textit{k3s}) dedicado, isolando cada usuário em um \textit{namespace} próprio e expondo os serviços via \textit{Ingress} ou \textit{NodePort}. Essa mudança de arquitetura, motivada pelas particularidades do provisionamento de contêineres, elimina a necessidade de um passo intermediário de configuração e aproxima a plataforma de práticas nativas de orquestração de contêineres, mantendo o objetivo original de reduzir a complexidade técnica exposta ao usuário final.




\par O código-fonte e o material complementar constam no Apêndice~\ref{apendice:material}.

\section{Objetivos}

\subsection*{Objetivo Geral}
Desenvolver e implementar uma Plataforma de Engenharia Self-Service capaz de automatizar e simplificar a provisão de ambientes de infraestrutura virtualizados \textit{on-premises} — máquinas virtuais e contêineres de aplicação —, abstraindo a complexidade técnica para o usuário final e acelerando o ciclo de pesquisa e desenvolvimento.

\subsection*{Objetivos Específicos}
\begin{itemize}
    \item Projetar e construir um portal web que permita aos usuários (alunos e pesquisadores) selecionar e personalizar as especificações de sua infraestrutura, como tipo de recurso (VM ou contêiner), quantidade de CPU, memória e armazenamento, dentro de limites definidos pelo perfil do usuário.
    \item Integrar o portal com o Terraform para orquestrar a criação e a remoção de máquinas virtuais no Proxmox VE, com base nas escolhas do usuário.
    \item Implementar o provisionamento de contêineres de aplicação em um cluster Kubernetes dedicado, incluindo isolamento por \textit{namespace}, aplicação de quotas por perfil e exposição de serviços via rede.
    \item Prover acesso interativo aos ambientes provisionados, por meio de terminal web tanto para máquinas virtuais (via proxy WebSocket para o Proxmox) quanto para contêineres (via \textit{exec} no cluster Kubernetes).
    \item Avaliar a eficácia da plataforma em termos de redução do tempo de provisão de infraestrutura, diminuição da dependência de equipes de TI para configurações básicas e melhoria da experiência dos pesquisadores no acesso a ambientes dedicados para suas atividades de desenvolvimento e pesquisa.
\end{itemize}


\clearpage

	\chapter{Fundamentação teórica}
\label{chapter:fundamentacao_teorica}

Para uma melhor compreensão da implementação e da metodologia, é necessário explicar os principais conceitos e ferramentas que fundamentam o desenvolvimento da plataforma proposta.

%--------------------------------Seção 2.1--------------------------------%

\section{Infraestrutura como Código e Provisionamento Automatizado}

\par Infraestrutura como Código (\textit{Infrastructure as Code} — IaC) é a prática de gerenciar e provisionar infraestrutura computacional por meio de arquivos de configuração legíveis por máquina, em vez de processos manuais e interativos. Essa abordagem permite versionar, revisar e reproduzir ambientes de forma consistente, reduzindo a variação entre ambientes e a incidência de erros humanos \cite{KunalBhandari}.\\
\par Estudos sobre automação de infraestrutura demonstram que o provisionamento manual de ambientes computacionais é um processo demorado e sujeito a falhas, sendo a automação de ponta a ponta um fator determinante para a redução do tempo de entrega de recursos \cite{Sandobalín}. De forma semelhante, abordagens baseadas em modelos para o provisionamento de \textit{hardware} têm sido propostas como forma de reduzir a intervenção manual e padronizar a criação de ambientes de TI \cite{Gokhale}.

\subsection{Terraform}
O Terraform é uma ferramenta de código aberto de Infraestrutura como Código, mantida pela HashiCorp, que permite declarar recursos de infraestrutura (como máquinas virtuais, redes e discos) em arquivos de configuração e aplicá-los de forma idempotente contra provedores de nuvem ou de virtualização \textit{on-premises} \cite{terraformDocs:online}. O Terraform mantém um arquivo de estado que reflete os recursos efetivamente provisionados, permitindo identificar divergências entre a configuração declarada e o estado real da infraestrutura, e aplicar apenas as mudanças necessárias a cada execução.

%------------------------------- Seção 2.2 --------------------------------------------%

\section{Engenharia de Plataforma e Infraestrutura Self-Service}

\par A Engenharia de Plataforma (\textit{Platform Engineering}) é a disciplina que se concentra em projetar e construir ferramentas, portais e fluxos de trabalho que reduzam a carga cognitiva de desenvolvedores e pesquisadores no consumo de infraestrutura, oferecendo capacidades de autoatendimento (\textit{self-service}) sobre uma base de automação subjacente \cite{Sandobalín}. O objetivo central dessa abordagem é permitir que o usuário final solicite e obtenha um ambiente pronto para uso sem depender diretamente de uma equipe de operações para cada solicitação.\\
\par Pesquisas sobre o impacto do desenvolvimento de infraestrutura na produtividade indicam que a disponibilidade de recursos computacionais sob demanda, com baixo tempo de espera, está diretamente associada a ganhos de eficiência em ambientes de pesquisa e desenvolvimento \cite{Muhammad}. Nesse sentido, plataformas self-service buscam otimizar a chamada Experiência do Pesquisador (\textit{Researcher Experience} — RX), termo que descreve a facilidade com que um pesquisador consegue obter e utilizar os recursos computacionais necessários às suas atividades, sem se envolver com os detalhes de baixo nível da infraestrutura subjacente \cite{Minoru}.

\subsection{Modelos de recursos: máquinas virtuais e contêineres}
Uma plataforma self-service de infraestrutura pode expor ao usuário diferentes modelos de recurso computacional. Máquinas virtuais oferecem isolamento completo em nível de sistema operacional, sendo adequadas a cargas de trabalho que exigem um ambiente completo e persistente. Contêineres de aplicação, por outro lado, compartilham o núcleo do sistema operacional do hospedeiro e são orquestrados por um sistema de gerenciamento de cargas de trabalho, sendo mais adequados a aplicações efêmeras ou que já seguem um modelo de implantação baseado em imagens.

%----------------------------------------- Seção 2.3 ------------------------------------------------

\section{Virtualização e Orquestração de Contêineres}

\subsection{Proxmox VE}
O Proxmox VE (\textit{Virtual Environment}) é uma plataforma de virtualização de código aberto que combina KVM (\textit{Kernel-based Virtual Machine}) para máquinas virtuais e LXC para contêineres do sistema operacional, oferecendo uma API REST completa para automação de todas as operações de gerenciamento, como criação, consulta de status e acesso a terminal das VMs \cite{proxmoxDocs:online}. Por expor uma API REST e suportar autenticação via \textit{token}, o Proxmox VE viabiliza a integração com ferramentas de automação externas, como o Terraform, permitindo que a criação de máquinas virtuais seja declarada como código em vez de realizada manualmente pela interface web.

\subsection{Kubernetes e k3s}
O Kubernetes é uma plataforma de orquestração de contêineres de código aberto que automatiza a implantação, o escalonamento e o gerenciamento de aplicações em contêineres \cite{kubernetesDocs:online}. Entre seus principais conceitos estão o \textit{namespace}, que permite isolar logicamente recursos de diferentes usuários ou equipes dentro do mesmo cluster; o \textit{Ingress}, responsável por rotear tráfego HTTP externo até os serviços internos com base em regras de domínio; e o \textit{NodePort}, que expõe um serviço em uma porta fixa de cada nó do cluster, sendo útil para protocolos não HTTP.\\
\par O k3s é uma distribuição leve do Kubernetes, empacotada em um único binário, voltada para ambientes com recursos computacionais mais restritos ou para cenários de borda (\textit{edge}), mantendo compatibilidade com a API padrão do Kubernetes e permitindo o uso das mesmas ferramentas de administração do ecossistema.

%----------------------------------------- Seção 2.4 ------------------------------------------------

\section{Autenticação e Autorização}

\par O controle de acesso de uma plataforma multiusuário exige um mecanismo que permita identificar o usuário autenticado em cada requisição, sem a necessidade de manter estado de sessão no servidor. O JSON Web Token (JWT) é um padrão aberto (RFC 7519) que define um formato compacto e autocontido para transmitir informações de forma segura entre partes, como um conjunto de reivindicações (\textit{claims}) assinadas digitalmente \cite{jwtRfc:online}. Por ser autocontido e verificável apenas com a chave de assinatura, o JWT permite que um serviço \textit{backend} valide a identidade e o perfil de um usuário em cada requisição sem consultar uma base de sessões centralizada, sendo amplamente utilizado em arquiteturas de API REST.

	\chapter{Metodologia}
\label{chapter:metodologia}

Este capítulo descreve, em alto nível, como o trabalho foi conduzido: a classificação da pesquisa, a elicitação dos requisitos, as etapas de desenvolvimento da solução e a forma de validação. Os detalhes técnicos do que foi construído são apresentados no Capítulo~\ref{chapter:implementacao}.

\section{Caracterização da pesquisa}

Este trabalho adotou a abordagem de Prova de Conceito (PoC), em um estudo de caso aplicado ao ambiente de virtualização \textit{on-premises} de um laboratório acadêmico, para validar a viabilidade técnica e prática de uma Plataforma de Engenharia Self-Service (\textit{Internal Research Platform}) voltada à automação da infraestrutura virtualizada. O objetivo geral é o descrito na Seção~1.1: automatizar e simplificar o provisionamento de máquinas virtuais e contêineres de aplicação, abstraindo a complexidade técnica para o usuário final.\\
\par A PoC foi conduzida de maneira incremental, com o sistema evoluindo em fases sucessivas, cada uma agregando uma capacidade completa e testável à plataforma. A escolha se justifica por tratar-se de uma solução que integra múltiplas ferramentas de infraestrutura (Proxmox, Terraform, Kubernetes) em um único portal \textit{self-service}, cujo objetivo principal é validar a viabilidade e os benefícios dessa integração de forma prática e demonstrável. O desenvolvimento incremental permitiu identificar rapidamente obstáculos técnicos — como a necessidade de abandonar o Ansible em favor de uma orquestração nativa em Kubernetes para contêineres — e ajustar a arquitetura ao longo do trabalho.

\section{Elicitação de requisitos}
\label{sec:elicitacao}

Os requisitos foram levantados a partir do problema observado no ambiente do laboratório — a solicitação de máquinas virtuais depende de interação direta com quem administra o Proxmox VE — e das definições do pré-projeto e das reuniões de orientação. Definiram-se dois perfis de usuário, com limites de recursos distintos: o \textit{aluno}, com limites menores, e o \textit{pesquisador}, com limites maiores e responsável pelo cadastro de novos usuários. A Tabela~\ref{tab:requisitos} resume os principais requisitos identificados. O levantamento partiu da demanda apresentada pelo professor Elton Bezerra Torres, orientador deste trabalho, que já dispunha da base dos requisitos que orientaram o projeto; toda a informação inicial passou por ele, por meio de áudios, documentos e reuniões de requisitos realizadas por videoconferência.

\begin{table}[htb]
    \centering
    \caption{Principais requisitos da plataforma}
    \label{tab:requisitos}
    \begin{tabularx}{\textwidth}{|l|X|}
        \hline
        \textbf{Id} & \textbf{Requisito} \\ \hline
        \multicolumn{2}{|l|}{\textit{Requisitos funcionais}} \\ \hline
        RF01 & Autenticar usuários e diferenciar os perfis aluno e pesquisador \\ \hline
        RF02 & Cadastro de usuários realizado apenas por pesquisadores (sem autocadastro) \\ \hline
        RF03 & Criar, listar, ligar, desligar e excluir máquinas virtuais, dentro dos limites do perfil \\ \hline
        RF04 & Criar, listar e excluir contêineres de aplicação, dentro das quotas do perfil \\ \hline
        RF05 & Acessar o terminal das VMs e dos contêineres pelo navegador \\ \hline
        RF06 & Expor serviços dos contêineres na rede (HTTP ou TCP) \\ \hline
        RF07 & Permitir o uso de GPU em cargas de IA/ML, anexando uma GPU física à máquina virtual (\textit{pass-through}) \\ \hline
        \multicolumn{2}{|l|}{\textit{Requisitos não funcionais}} \\ \hline
        RNF01 & Provisionamento reprodutível, definido como código \\ \hline
        RNF02 & Credenciais da infraestrutura nunca expostas ao usuário final \\ \hline
        RNF03 & Isolamento entre usuários \\ \hline
        RNF04 & Interface simples, sem exigir conhecimento da infraestrutura subjacente \\ \hline
    \end{tabularx}
    \par\small Fonte: elaborado pelo autor.
\end{table}

\par Como o desenvolvimento foi incremental, esse levantamento inicial foi refinado ao longo das fases, à medida que cada capacidade era detalhada nas reuniões de orientação. Ao final, os requisitos foram consolidados e rastreados até a implementação, tomando como base as rotas da API, com seus métodos HTTP, e as regras de negócio aplicadas no \textit{backend}. Todos os requisitos da Tabela~\ref{tab:requisitos} têm correspondência na implementação. Além deles, o refinamento resultou nos requisitos da Tabela~\ref{tab:requisitos_complementares}, que complementam os anteriores sem substituí-los. A coluna \textit{Evidência} indica a rota ou a regra do sistema de onde cada requisito foi derivado.

{\small
\begin{longtable}{|l|p{7.4cm}|p{4.9cm}|}
    \caption{Requisitos complementares derivados do sistema implementado}
    \label{tab:requisitos_complementares} \\
    \hline
    \textbf{Id} & \textbf{Requisito} & \textbf{Evidência} \\ \hline
    \endfirsthead
    \hline
    \textbf{Id} & \textbf{Requisito} & \textbf{Evidência} \\ \hline
    \endhead
    \multicolumn{3}{|l|}{\textit{Requisitos funcionais}} \\ \hline
    RF08 & Consultar os detalhes de uma máquina virtual e reiniciá-la & \path{GET /api/vms/:id}; \path{POST /api/vms/:id/reboot} \\ \hline
    RF09 & Gerar, para VMs Ubuntu, usuário e senha de acesso por VM, exibidos ao usuário uma única vez, na criação & Resposta \path{201} de \path{POST /api/vms} (campo \textit{credentials}) \\ \hline
    RF10 & Impedir que uma VM com GPU seja ligada enquanto a placa estiver em uso por outra VM, informando o conflito ao usuário & \path{POST /api/vms/:id/start} (resposta \path{409}) \\ \hline
    RF11 & Consultar o catálogo de imagens disponíveis para contêineres & \path{GET /api/containers/catalog} \\ \hline
    RF12 & Configurar, na criação do contêiner, variáveis de ambiente, comando, porta e protocolo (HTTP ou TCP) e volume persistente & Corpo de \path{POST /api/containers} \\ \hline
    RF13 & Consultar os detalhes e os \textit{logs} de um contêiner & \path{GET /api/containers/:id}; \path{GET /api/containers/:id/logs} \\ \hline
    RF14 & Consultar e atualizar os dados do próprio perfil e alterar a própria senha, mediante a senha atual & \path{GET} e \path{PATCH /api/me}; \path{PUT /api/me/password} \\ \hline
    RF15 & Listar os usuários cadastrados, com filtro por perfil, alterar o \textit{e-mail} e excluir usuários (apenas pesquisador), removendo também as VMs e os contêineres do usuário excluído & \path{GET /api/users}; \path{PATCH} e \path{DELETE /api/users/:id} \\ \hline
    \multicolumn{3}{|l|}{\textit{Requisitos não funcionais}} \\ \hline
    RNF05 & Limites de recursos por perfil: aluno com até 2 CPUs, 4 GB de memória, 5 contêineres e 20 GB de armazenamento persistente; pesquisador com até 8 CPUs, 16 GB, 20 contêineres e 100 GB. Perfil desconhecido recebe os limites do aluno & Regras de limite do \textit{backend}; resposta de erro de limite excedido \\ \hline
    RNF06 & Todas as rotas, exceto o \textit{login}, exigem autenticação por \textit{token} JWT com validade configurável; a gestão de usuários é restrita ao perfil pesquisador & \textit{Middleware} de autenticação e de perfil; respostas \path{401} e \path{403} \\ \hline
    RNF07 & Senhas armazenadas apenas na forma de \textit{hash} (\textit{bcrypt}), com tamanho mínimo de 8 caracteres & Cadastro e troca de senha \\ \hline
    RNF08 & Cada usuário acessa somente as suas próprias VMs e contêineres; contêineres ficam em \textit{namespace} próprio, com quota; o \textit{resource pool} das VMs é definido pelo \textit{backend} conforme o perfil, nunca pelo usuário & Verificação de propriedade nos casos de uso; \textit{namespace} por usuário \\ \hline
    RNF09 & Erros internos da infraestrutura (Proxmox VE, Terraform, Kubernetes) não são repassados ao usuário, que recebe mensagens genéricas em português & Tratamento de erros nos \textit{handlers} \\ \hline
    RNF10 & Configuração e segredos fornecidos por variáveis de ambiente, sem valores sensíveis no código-fonte & Módulo de configuração do \textit{backend} \\ \hline
    RNF11 & API documentada no padrão OpenAPI (Swagger) & Rota \path{/swagger} e arquivo \path{swagger.yaml} \\ \hline
    RNF12 & Código organizado em camadas desacopladas por interfaces, com testes automatizados e testes de fumaça contra o ambiente real & Estrutura do \textit{backend}; \path{cmd/proxmox-check} \\ \hline
\end{longtable}
}

\section{Desenvolvimento da solução}
\label{sec:desenvolvimento}

O desenvolvimento foi dividido nas seguintes fases, cada uma entregando uma capacidade completa e validada antes do início da seguinte:

\begin{enumerate}
    \item \textbf{Autenticação e usuários}: mecanismo de autenticação e perfis de acesso;
    \item \textbf{Máquinas virtuais no Proxmox}: integração inicial com a API do Proxmox VE para criação, consulta e remoção de VMs;
    \item \textbf{Terminal web}: acesso interativo às VMs diretamente pelo navegador;
    \item \textbf{Integração \textit{frontend}/\textit{backend}}: portal web conectado ao \textit{backend}, cobrindo o fluxo principal para o usuário final;
    \item \textbf{Migração para Infraestrutura como Código}: substituição do provisionamento imperativo por declarações Terraform;
    \item \textbf{Contêineres de aplicação em k3s}: segundo modelo de recurso, sobre um \textit{cluster} Kubernetes dedicado;
    \item \textbf{Gestão de usuários e perfil}: cadastro por pesquisadores, tela de perfil e listagem/exclusão de usuários.
\end{enumerate}

\par Estas fases podem ser vistas na Figura~\ref{fig:fases_desenvolvimento}.

\figimg{fig:fases_desenvolvimento}{Etapas de desenvolvimento da PoC}{fases_desenvolvimento}{elaboração própria}

\section{Validação}

A validação consistiu em testes de fluxo de trabalho executados contra o ambiente real do laboratório (Proxmox VE e \textit{cluster} k3s), tanto por utilitários de linha de comando quanto pelo uso ponta a ponta do portal, conforme descrito no Capítulo~\ref{chapter:experimentos}. Os resultados são confrontados com os objetivos específicos no Capítulo~\ref{chapter:resultados}.

	\chapter{Implementação}
\label{chapter:implementacao}

O objetivo deste capítulo é descrever como a Plataforma de Engenharia Self-Service — ProxLab — foi implementada, detalhando a arquitetura do \textit{backend}, os fluxos de provisionamento de máquinas virtuais e de contêineres, e os mecanismos de autenticação e acesso remoto oferecidos ao usuário final, conforme a visão geral da Figura~\ref{fig:arquitetura_proxlab}.

\figimg{fig:arquitetura_proxlab}{Arquitetura geral do ProxLab}{arquitetura_proxlab}{elaboração própria}

\section{Arquitetura em camadas}

O \textit{backend} foi desenvolvido em Go, utilizando o \textit{framework} web Gin e o ORM GORM sobre um banco de dados PostgreSQL. A aplicação segue uma arquitetura em camadas, organizada de forma que cada camada dependa apenas de interfaces, e não de implementações concretas (Figura~\ref{fig:camadas_backend}):

\begin{itemize}
    \item \texttt{handler} — controladores HTTP (Gin), responsáveis por validar o \textit{payload} da requisição e delegar a lógica de negócio;
    \item \texttt{usecase} — lógica de negócio da aplicação, independente de detalhes de infraestrutura;
    \item \texttt{domain/repository} — interfaces de persistência, nunca implementações;
    \item \texttt{repository} — implementações concretas das interfaces de persistência, utilizando GORM.
\end{itemize}

Essa separação é reforçada por uma regra de desenvolvimento adotada no projeto: nenhuma lógica de negócio deve residir no ponto de entrada da aplicação, e nenhuma chamada ao GORM deve ocorrer diretamente em um \textit{handler}.

\figimg{fig:camadas_backend}{Arquitetura em camadas do \textit{backend}}{camadas_backend}{elaboração própria}

\section{Provisionamento de máquinas virtuais}

O fluxo de criação de uma máquina virtual segue o caminho \textit{Handler} \(\rightarrow\) \textit{UseCase} \(\rightarrow\) \textit{Provisioner} \(\rightarrow\) \texttt{terraform-exec} \(\rightarrow\) Proxmox VE. Ao receber uma solicitação de criação de VM, o \textit{backend} gera um arquivo de configuração Terraform (\texttt{main.tf}) em um \textit{workspace} isolado por identificador único de VM, utilizando o \textit{provider} \texttt{bpg/proxmox}, e executa o Terraform por meio da biblioteca \texttt{terraform-exec}. As credenciais de acesso ao Proxmox são fornecidas apenas via variáveis de ambiente do processo, nunca gravadas em arquivo de configuração.\\

\par Máquinas virtuais Ubuntu são criadas a partir da clonagem completa de um \textit{template} \textit{cloud-init} previamente preparado no Proxmox, já nascendo ligadas e com um usuário padrão cuja senha é gerada de forma aleatória por VM e devolvida uma única vez na resposta da criação, nunca persistida em banco de dados. Máquinas virtuais de outros sistemas operacionais seguem o fluxo de VM vazia, cabendo ao usuário realizar a instalação do sistema operacional.\\

\par Quando o usuário solicita GPU, o \textit{backend} escolhe um \textit{resource mapping} PCI do Proxmox VE (o que tiver menos máquinas associadas) e o anexa à VM por \textit{pass-through}, por meio de um bloco \texttt{hostpci} na configuração Terraform. Como a placa não pode ser compartilhada, a VM com GPU nasce desligada e a operação de ligar é serializada: antes de iniciá-la, o \textit{backend} verifica se outra VM está usando o mesmo \textit{mapping} e, se o Proxmox VE recusar a inicialização, a falha é traduzida em uma resposta de conflito, sem expor ao usuário a mensagem interna da infraestrutura. A instalação automática do \textit{driver} de vídeo e do Docker nessas VMs, por \textit{cloud-init}, é descrita no Apêndice~\ref{apendice:material}.\\

\par Operações de consulta de status, ligar e desligar uma VM já existente não passam pelo Terraform: são realizadas diretamente contra a API REST do Proxmox VE, por meio de um cliente HTTP dedicado (\textit{ProxmoxClient}), autenticado via \textit{token} de API. A Figura~\ref{fig:tela_criar_vm} apresenta a tela de criação de máquina virtual utilizando o portal do ProxLab.

\figimg{fig:tela_criar_vm}{Tela de criação de máquina virtual no portal}{front-tela-criar-vm}{captura de tela do ProxLab}

\section{Provisionamento de contêineres de aplicação}

Contêineres de aplicação seguem um fluxo distinto: \textit{Handler} \(\rightarrow\) \textit{UseCase} \(\rightarrow\) \textit{ContainerRuntime} \(\rightarrow\) \texttt{client-go} \(\rightarrow\) cluster k3s. O cluster é executado em uma máquina virtual dedicada, provisionada uma única vez por meio de um utilitário de \textit{bootstrap} que instala o k3s e gera o \textit{kubeconfig} utilizado pelo \textit{backend}.\\

\par Cada usuário da plataforma possui um \textit{namespace} próprio no cluster, identificado por um sufixo derivado do seu identificador, e uma quota de recursos (CPU, memória) definida de acordo com o seu perfil (aluno ou pesquisador). O estado de um contêiner não é persistido em banco de dados: ele é lido diretamente do cluster a cada requisição, garantindo que a plataforma sempre reflita o estado real do recurso.\\

\par A exposição de rede de um contêiner depende do tipo de acesso solicitado pelo usuário: quando o contêiner expõe uma porta HTTP e um domínio de \textit{Ingress} está configurado, o acesso é feito por meio de um subdomínio dedicado; quando o acesso é via TCP, ou quando não há domínio configurado, o serviço é exposto por meio de \textit{NodePort}, em uma porta na faixa 30000–32767 do nó do cluster. Um contêiner também pode ser criado apenas para acesso via terminal, sem nenhuma porta de rede exposta. A Figura~\ref{fig:tela_criar_container} apresenta a tela de criação de contêineres usando o portal do ProxLab.

\figimg{fig:tela_criar_container}{Tela de criação de contêiner de aplicação no portal}{front-tela-criar-container}{captura de tela do ProxLab}

\section{Autenticação e gestão de usuários}

O acesso à plataforma é protegido por autenticação baseada em JWT. Após o \textit{login} (Figura~\ref{fig:tela_login}), o token é armazenado no \texttt{localStorage} do navegador e enviado em todas as requisições subsequentes por meio do cabeçalho \texttt{Authorization}. Conexões WebSocket, por não suportarem cabeçalhos HTTP customizados na negociação inicial, recebem o token por meio de um parâmetro de consulta (\textit{query param}) na URL de conexão. Não há autocadastro de usuários: o cadastro de novos usuários exige um token JWT de um usuário com perfil \textit{pesquisador}. O perfil define os limites de recursos e o \textit{pool} do Proxmox em que as VMs são criadas, sempre decidido pelo \textit{backend} e nunca pelo \textit{payload} da requisição. Os usuários podem alterar seu nome e senha em uma tela de perfil, e o pesquisador dispõe de uma tela de gestão que lista, filtra por perfil, altera o e-mail e exclui usuários (Figura~\ref{fig:tela_usuarios}); a exclusão remove antes todas as VMs e contêineres do usuário.

\figimg{fig:tela_login}{Tela de \textit{login} da plataforma}{front-tela-login(geral)}{captura de tela do ProxLab}

\figimg{fig:tela_usuarios}{Tela de usuários cadastrados (visão do pesquisador)}{front-tela-usuarios(pesquisador)}{captura de tela do ProxLab}

\section{Terminal web}

O acesso interativo aos recursos provisionados é oferecido diretamente pelo navegador, sem a necessidade de um cliente SSH local (Figura~\ref{fig:terminal_web}). Para máquinas virtuais, o \textit{backend} atua como um proxy WebSocket entre o navegador e o protocolo \textit{termproxy} do Proxmox VE, reproduzindo o seu \textit{handshake} específico: uma primeira mensagem de autenticação no formato \texttt{"<usuário>:<ticket>\textbackslash n"}, respondida com \texttt{OK}, seguida do enquadramento de dados de teclado e de redimensionamento de terminal segundo o protocolo do Proxmox. Para contêineres, o acesso a terminal é feito por meio da operação de \textit{exec} do próprio cluster Kubernetes, sem a necessidade de um protocolo intermediário.

\figimgh{fig:terminal_web}{Terminal web de uma máquina virtual Ubuntu no navegador}{front-tela-terminal-web-vm}{captura de tela do ProxLab}{0.17}

\section{Infraestrutura como código e configuração inicial}

A infraestrutura de máquinas virtuais é definida como código: cada VM possui seu próprio \textit{workspace} Terraform, de modo que criar ou destruir um recurso equivale a aplicar ou destruir uma configuração versionável e reprodutível. A configuração inicial das VMs Ubuntu (usuário, senha e chaves SSH) é entregue por \textit{cloud-init}, a partir do \textit{template} clonado. Não é utilizado o Ansible: a instalação de \textit{software} é oferecida por meio de contêineres de aplicação. A Figura~\ref{fig:terraform_main_tf} apresenta um exemplo de configuração usando o Terraform.

\figimgh{fig:terraform_main_tf}{Exemplo de configuração Terraform gerada pelo \textit{backend}}{terraform_main_tf}{elaboração própria (exemplo ilustrativo baseado no \textit{template} do \textit{backend})}{0.16}

	\chapter{Experimentos}
\label{chapter:experimentos}

Para validar o funcionamento da plataforma contra um ambiente Proxmox VE e um cluster k3s reais — e não apenas contra simulações —, foram conduzidas validações de fluxo de trabalho, executadas tanto por utilitários de linha de comando quanto pelo uso ponta a ponta do portal web. Cada validação percorre um fluxo completo, do ponto de vista do usuário, e verifica se o resultado esperado é observado no ambiente real.

\section{Validação de criação de recursos}

Foram implementados dois utilitários de linha de comando (\texttt{cmd/proxmox-check}) que exercitam o fluxo completo de provisionamento contra a infraestrutura real:

\begin{itemize}
    \item o primeiro cria uma máquina virtual por clone do \textit{template} Ubuntu com \textit{cloud-init} (o mesmo caminho utilizado pelo portal), consulta o seu status e a remove, validando a integração via Terraform e via API REST;
    \item o segundo, executado com a opção \texttt{--k3s}, cria, consulta e remove um contêiner de aplicação no cluster k3s real, validando a integração via \texttt{client-go}.
\end{itemize}

\par Em seguida, o mesmo fluxo foi executado pelo portal web, com a criação de uma máquina virtual Ubuntu — acessada via terminal web e removida — e de um contêiner de aplicação, conforme a Seção~\ref{sec:val_conteudo}. Foram verificados também os limites por perfil: a criação de VM por um usuário com perfil \textit{aluno} é feita no \textit{pool} definido pelo \textit{backend} para esse perfil, independentemente do valor enviado pelo cliente.

\figimg{fig:val_vm_proxmox}{Máquina virtual criada pela plataforma, vista na interface do Proxmox VE, com o histórico de tarefas (\textit{resize}, \textit{start} e a remoção da VM do teste de fumaça)}{val_vm_proxmox}{captura de tela do Proxmox VE}

\figimg{fig:val_saida_smoke}{Saída dos testes de fumaça contra o ambiente real: contêiner no cluster k3s (\texttt{--k3s}) e máquina virtual no Proxmox VE}{val_saida_smoke}{captura de tela do terminal}

\section{Validação de criação de conteúdo}
\label{sec:val_conteudo}

% TODO(autor): confirmar se "criação de conteúdo" no sentido do professor é o que foi feito dentro dos recursos (banco, dados, logs). Ajustar o texto se for outro entendimento.
Além de criar o recurso, validou-se que é possível gerar e utilizar conteúdo dentro dele. Foi criado, pelo portal, um contêiner de aplicação PostgreSQL com variáveis de ambiente, volume persistente e exposição de porta TCP via \textit{NodePort}; foram consultados os logs do contêiner e validada uma conexão via \texttt{psql}, na qual o banco respondeu às consultas, antes de o contêiner ser removido. De modo análogo, na VM Ubuntu, o terminal web foi utilizado para executar comandos no sistema criado.

\par O acesso externo a um contêiner por meio do par \texttt{ip:porta} do \textit{NodePort} não pôde ser validado durante o desenvolvimento, pois a rede VPN utilizada para acesso ao ambiente de laboratório bloqueia o tráfego de entrada nessas portas a partir de fora da rede interna. Essa validação depende de o \textit{backend} estar em execução na própria VM do laboratório, e é discutida como trabalho futuro no Capítulo~\ref{chapter:conclusoes_trabalhos_futuros}.

\figimg{fig:val_postgres}{Detalhes do contêiner PostgreSQL em execução, com endereço externo (\textit{NodePort}) e variáveis de ambiente}{front-tela-container-em-execucao}{captura de tela do ProxLab}

\figimg{fig:val_logs_container}{Terminal e logs do contêiner PostgreSQL: conexão via \texttt{psql} e listagem das tabelas criadas}{front-tela-terminal-container}{captura de tela do ProxLab}

\section{Validação de \textit{pass-through} de GPU}

% TODO(autor): GPU NÃO está implementada no código (fora da v1). Se o professor quer este experimento, é preciso
% implementar/testar o pass-through (PCIe passthrough no Proxmox / hostpci no Terraform) antes de preencher esta seção.
% Se não for possível até 15/11, manter como trabalho futuro e remover esta seção.
Esta validação tem por objetivo verificar se uma máquina virtual criada pela plataforma consegue acessar uma GPU física do servidor por meio de \textit{pass-through} de dispositivo PCIe, requisito de cargas de IA/ML. Até a presente versão, o acesso a GPU não foi implementado, razão pela qual não há resultados a reportar; o tema é tratado como trabalho futuro no Capítulo~\ref{chapter:conclusoes_trabalhos_futuros}.

\figplaceholder{fig:val_gpu}{GPU visível dentro da máquina virtual (\texttt{nvidia-smi})}{Captura do terminal da VM exibindo a saída de \texttt{nvidia-smi} (a ser feita somente se o pass-through for implementado)}{captura de tela do terminal}

\section{Validação de criação e autenticação de usuários}

Validou-se o fluxo de identidade da plataforma: (i) o \textit{login} retorna um \textit{token} JWT e as rotas protegidas rejeitam requisições sem \textit{token} ou com \textit{token} inválido (resposta 401, com redirecionamento para a tela de \textit{login}); (ii) o cadastro de um novo usuário só é aceito quando solicitado por um usuário com perfil \textit{pesquisador}, sendo recusado para o perfil \textit{aluno}; (iii) o usuário consegue alterar seu nome e sua senha, sendo recusada a troca quando a senha atual informada está incorreta; e (iv) o pesquisador consegue listar, alterar o e-mail e excluir usuários, sendo a exclusão acompanhada da remoção das VMs e dos contêineres do usuário.

\figimg{fig:val_usuarios}{Tela de cadastro de novo usuário pelo pesquisador}{front-tela-cadastro-usuario}{captura de tela do ProxLab}

	\chapter{Resultados}

\section{Resultados por objetivo}

Neste capítulo é feito um levantamento dos resultados obtidos frente aos objetivos específicos definidos na Seção~1.1, confrontando o que foi originalmente proposto com o que foi de fato implementado e validado.

\subsection{Portal web e seleção de recursos}
O portal web, desenvolvido em Next.js, permite ao usuário selecionar entre máquina virtual e contêiner de aplicação, respeitando os limites de CPU, memória e armazenamento definidos pelo seu perfil (aluno ou pesquisador). Esse objetivo foi integralmente atendido.

\subsection{Provisionamento de VMs via Terraform}
A integração do portal com o Terraform para a criação e a remoção de máquinas virtuais no Proxmox VE foi implementada e validada tanto por testes de fumaça quanto pelo uso ponta a ponta pelo portal, conforme descrito no Capítulo~4. Esse objetivo foi integralmente atendido.

\subsection{Automação de configuração: de Ansible a contêineres nativos}
O objetivo original de automação de configuração via Ansible — instalação de \textit{software}, bancos de dados e dependências dentro de uma VM genérica — não foi implementado na forma originalmente proposta. Em seu lugar, a plataforma passou a oferecer o provisionamento nativo de contêineres de aplicação em um cluster Kubernetes, o que resolve o mesmo problema de fundo (disponibilizar \textit{software} pronto para uso sem intervenção manual do usuário) por um caminho arquiteturalmente mais próximo das práticas atuais de orquestração de contêineres. Considera-se esse objetivo atendido por uma via alternativa à originalmente prevista.

\subsection{Acesso interativo aos ambientes}
O acesso via terminal web foi implementado tanto para máquinas virtuais (proxy WebSocket para o protocolo \textit{termproxy} do Proxmox) quanto para contêineres (\textit{exec} via cluster Kubernetes), superando o escopo originalmente previsto no pré-projeto, que não especificava esse mecanismo de acesso.

\subsection{Redução de tempo, autonomia e padronização}
Observou-se, ao longo do uso da plataforma nos experimentos do Capítulo~5, que o provisionamento de um ambiente — VM ou contêiner — ocorre de forma autônoma pelo usuário, sem intervenção de uma equipe de infraestrutura, e em poucos minutos, contra um processo manual que exigiria a criação e configuração individual de cada recurso diretamente no Proxmox ou no cluster. O uso de Terraform como fonte declarativa da infraestrutura de VMs também contribui para a padronização e a reprodutibilidade dos ambientes criados.

\subsection{Acesso a GPU}
O acesso a GPU para cargas de trabalho de IA/ML, previsto como objetivo específico no pré-projeto, não foi implementado até o presente momento, permanecendo como trabalho futuro, conforme discutido no Capítulo~\ref{chapter:conclusoes_trabalhos_futuros}.

	\chapter{Conclusões e Trabalhos Futuros}
\label{chapter:conclusoes_trabalhos_futuros}

\section{Conclusões}

Com base nos objetivos estabelecidos, pode-se afirmar que a Plataforma de Engenharia Self-Service desenvolvida neste trabalho atingiu seu objetivo geral: oferecer a alunos e pesquisadores um portal capaz de provisionar máquinas virtuais e contêineres de aplicação de forma autônoma, sem exigir conhecimento aprofundado da infraestrutura subjacente. A arquitetura final, sustentada por Terraform para VMs e por um cluster Kubernetes dedicado para contêineres, diverge da abordagem originalmente prevista no pré-projeto — que previa o uso de Ansible para configuração de ambientes —, mas preserva a ideia central de reduzir a complexidade técnica exposta ao usuário final, tendo se mostrado, na prática, uma solução mais aderente ao modelo de provisionamento de contêineres adotado.

\par O acesso a GPU para cargas de IA e \textit{machine learning} também foi atendido: o \textit{backend} anexa à máquina virtual, por \textit{pass-through}, uma das duas placas do servidor, balanceando o uso entre elas e impedindo que duas máquinas liguem a mesma placa ao mesmo tempo. A validação em ambiente real (Seção~\ref{sec:val_gpu}) confirmou a GPU visível dentro da VM.

\par Como limitações, destaca-se que a validação foi funcional, conduzida pelo autor em um único nó Proxmox VE e em um cluster k3s de nó único, sem medição controlada de tempo nem participação de um grupo de usuários. Ainda assim, os fluxos completos de criação, uso e remoção de recursos foram exercitados contra a infraestrutura real do laboratório, o que sustenta a viabilidade técnica da proposta.

\section{Trabalhos Futuros}
\label{sec:trabalhos_futuros}

Entre os itens que permanecem como trabalho futuro, destacam-se:

\begin{itemize}
    \item Ciclo de vida mais completo para contêineres de aplicação (\textit{start}, \textit{stop}, \textit{restart}) e suporte a \textit{build} de imagens a partir de \textit{Dockerfile} ou \textit{tarball} enviado pelo usuário.

    \item TLS no \textit{Ingress} dos contêineres expostos via domínio, hoje disponibilizados apenas em HTTP.

    \item Validação do acesso externo a contêineres por meio do par \texttt{ip:porta} do \textit{NodePort}, não realizada durante o desenvolvimento devido a restrições da rede VPN utilizada para acesso ao ambiente de laboratório.

    \item Avaliação quantitativa da plataforma, com a medição controlada do tempo de provisionamento em comparação ao processo manual e a aplicação de um questionário de usabilidade a alunos e pesquisadores do laboratório.

    \item Conclusão do cadastro por convite, em que o aluno ativa a própria conta no primeiro acesso por meio de um código enviado por e-mail, e integração ao \textit{backend} da instalação automática de drivers em VMs com GPU (Apêndice~A).

\end{itemize}

Através dessas frentes, a plataforma pode evoluir de uma Prova de Conceito validada para uma ferramenta mais completa de apoio à pesquisa e ao ensino no contexto do laboratório.
\clearpage
	
	
	\fancyhead[RO]{\thepage}
    \printbibliography[heading=bibintoc]
    \clearpage
    
    \fancyhead[LO]{\MakeUppercase{\appendixname}\,\thechapter}
	\begin{appendix}
	    \chapter{Apêndice}
\label{apendice:material}

Este apêndice reúne material complementar à implementação da plataforma.

\section{Repositório do código-fonte}

O código-fonte do ProxLab (\textit{backend}, \textit{frontend}, migrações e documentação de projeto) está disponível em:

\begin{center}
    \url{https://github.com/Brazf/platform-engineering}
\end{center}

\par O código é distribuído sob a licença Apache, versão 2.0, cujo texto integral acompanha o repositório no arquivo \texttt{LICENSE}.

\section{Variáveis de ambiente do \textit{backend}}

A configuração do \textit{backend} é feita exclusivamente por variáveis de ambiente. Os valores secretos (senhas, \textit{token} do Proxmox VE e segredo do JWT) nunca são versionados.

\begin{longtable}{|p{6.2cm}|p{8.4cm}|}
    \caption{Variáveis de ambiente do \textit{backend}}
    \label{tab:variaveis_ambiente} \\
    \hline
    \textbf{Variável} & \textbf{Descrição (valor padrão)} \\ \hline
    \endfirsthead
    \hline
    \textbf{Variável} & \textbf{Descrição (valor padrão)} \\ \hline
    \endhead
    \texttt{APP\_ENV} & Ambiente de execução; em desenvolvimento executa a migração automática do banco (\texttt{development}) \\ \hline
    \texttt{SERVER\_PORT} & Porta HTTP do servidor (\texttt{8080}) \\ \hline
    \texttt{DB\_HOST}, \texttt{DB\_PORT}, \texttt{DB\_USER}, \texttt{DB\_PASSWORD}, \texttt{DB\_NAME} & Conexão com o PostgreSQL \\ \hline
    \texttt{JWT\_SECRET} & Segredo de assinatura dos \textit{tokens} JWT (obrigatória) \\ \hline
    \texttt{JWT\_EXPIRY\_HOURS} & Validade do \textit{token} em horas (\texttt{24}) \\ \hline
    \texttt{PROXMOX\_HOST} & URL da API do Proxmox VE (obrigatória) \\ \hline
    \texttt{PROXMOX\_NODE} & Nó do Proxmox onde os recursos são criados \\ \hline
    \texttt{PROXMOX\_TOKEN\_ID}, \texttt{PROXMOX\_TOKEN\_SECRET} & Identificador e segredo do \textit{token} de API (obrigatórias) \\ \hline
    \texttt{PROXMOX\_POOL\_ALUNO}, \texttt{PROXMOX\_POOL\_PESQUISADOR} & \textit{Resource pool} das VMs de cada perfil, definido pelo \textit{backend} (\texttt{researchers}) \\ \hline
    \texttt{PROXMOX\_UBUNTU\_TEMPLATE\_VMID} & VMID do \textit{template} \textit{cloud-init} Ubuntu clonado nas VMs Ubuntu (\texttt{9000}) \\ \hline
    \texttt{PROXMOX\_GPU\_MAPPINGS} & Nomes dos \textit{resource mappings} PCI das GPUs, separados por vírgula; a ordem desempata a escolha; vazio desabilita a GPU \\ \hline
    \texttt{TERRAFORM\_BIN} & Caminho do binário do Terraform (\texttt{terraform}) \\ \hline
    \texttt{TERRAFORM\_WORKDIR} & Raiz dos \textit{workspaces} Terraform, um por VM \\ \hline
    \texttt{K3S\_NODE\_IP}, \texttt{K3S\_NODE\_GATEWAY} & Endereço fixo (CIDR) e \textit{gateway} da VM do cluster k3s \\ \hline
    \texttt{K3S\_NODE\_POOL} & \textit{Resource pool} da VM do cluster (\texttt{researchers}) \\ \hline
    \texttt{K3S\_NODE\_CPU}, \texttt{K3S\_NODE\_MEMORY\_MB}, \texttt{K3S\_NODE\_DISK\_GB} & Recursos da VM do cluster (\texttt{8}, \texttt{16384}, \texttt{200}) \\ \hline
    \texttt{INGRESS\_DOMAIN} & Domínio dos contêineres HTTP; vazio expõe por \textit{NodePort} \\ \hline
    \texttt{KUBECONFIG\_PATH} & Arquivo de acesso ao cluster; ausente, as rotas de contêineres respondem 503 \\ \hline
\end{longtable}

\section{Exemplo de configuração Terraform gerada com GPU}

O trecho a seguir mostra, de forma resumida, o \texttt{main.tf} gerado para uma VM Ubuntu com GPU (reticências indicam blocos omitidos). O bloco \texttt{hostpci} referencia o \textit{resource mapping} escolhido pelo \textit{backend}, e a VM nasce desligada (\texttt{started = false}).

\medskip
\noindent\begin{minipage}{\textwidth}
\footnotesize
\begin{verbatim}
resource "proxmox_virtual_environment_vm" "vm" {
  node_name     = "pve01"
  name          = "minha-vm"
  pool_id       = "researchers"
  started       = false

  clone {
    vm_id = 9000
    full  = true
  }

  cpu    { cores = 2 }
  memory { dedicated = 2048 }
  # disk, network_device, initialization, agent ...

  hostpci {
    device  = "hostpci0"
    mapping = "gpu-rtx4060-01"
    pcie    = true
  }
}
\end{verbatim}
\end{minipage}
\medskip

\section{Permissão do \textit{token} sobre as GPUs}

Para anexar uma GPU, o \textit{token} de API do ProxLab precisa da permissão \texttt{Mapping.Use} em cada \textit{resource mapping}. Ela é concedida pelo administrador do Proxmox VE, uma vez por \textit{mapping}, com o comando:

\medskip
\noindent\begin{minipage}{\textwidth}
\footnotesize
\begin{verbatim}
pveum acl modify /mapping/pci/gpu-rtx4060-01 \
    --tokens '<id-do-token>' --roles PVEMappingUser
\end{verbatim}
\end{minipage}
\medskip

Sem essa permissão, a criação de VM com GPU falha e o usuário recebe a mensagem de que a GPU está temporariamente indisponível. Além disso, a opção \texttt{pcie = true} exige que o \textit{template} usado na clonagem seja da máquina do tipo \texttt{q35}.

\section{\textit{Drivers} e Docker em VMs com GPU}

Anexar a GPU à VM não instala o \textit{driver} de vídeo no sistema convidado. Para abstrair essa dificuldade do usuário, a infraestrutura do laboratório adota o \textit{cloud-init}, com um arquivo \textit{vendor-data} preparado pelo administrador. Esta seção documenta o procedimento; os arquivos de referência (\path{create-template.sh}, \path{vendor-data-gpu.yaml} e o exemplo Terraform \texttt{04-vm-cloud-init}) estão no repositório de infraestrutura do laboratório\footnote{\url{https://github.com/ifpebj-ti/proxmox-infra-cces}}. Nesta versão do ProxLab, o \textit{backend} cria a VM e anexa a GPU; a instalação do \textit{driver} descrita a seguir é o procedimento documentado pela infraestrutura e não está integrada ao \textit{backend}.

Os dois arquivos de referência desta seção estão publicados nos endereços abaixo:

\par Exemplo Terraform de VM com \textit{cloud-init}:
{\small\begin{center}
\url{https://github.com/ifpebj-ti/proxmox-infra-cces/tree/main/terraform/examples/04-vm-cloud-init}
\end{center}}

\par \textit{Vendor-data} de GPU:
{\small\begin{center}
\url{https://github.com/ifpebj-ti/proxmox-infra-cces/blob/main/cloud-init/snippets/vendor-data-gpu.yaml}
\end{center}}

\subsection{Etapas}

\begin{enumerate}
    \item \textbf{Template (administrador, uma vez).} Criar no host Proxmox VE o \textit{template} \textit{cloud-init} Ubuntu 24.04 (nome \texttt{ubuntu-2404-cloud}), com o \textit{script} de criação de \textit{template} do repositório de infraestrutura.
    \item \textbf{\textit{Vendor-data} (administrador, uma vez).} Copiar o arquivo de \textit{vendor-data} de GPU para o diretório de \textit{snippets} do \textit{storage} do Proxmox VE, o que o torna disponível pelo identificador mostrado abaixo.
    \item \textbf{VM (a cada criação).} Na configuração Terraform da VM, além do bloco \texttt{hostpci} da seção anterior, referenciar o \textit{snippet} no bloco \texttt{initialization}.
\end{enumerate}

\medskip
\noindent\begin{minipage}{\textwidth}
\footnotesize
\begin{verbatim}
# Administrador, no host Proxmox VE (uma vez):
scp vendor-data-gpu.yaml \
    root@<host>:/datapool/templates/snippets/
# Identificador do snippet para o Terraform:
#   datapool-templates:snippets/vendor-data-gpu.yaml
\end{verbatim}
\end{minipage}
\medskip

\medskip
\noindent\begin{minipage}{\textwidth}
\footnotesize
\begin{verbatim}
initialization {
  vendor_data_file_id = "datapool-templates:snippets/vendor-data-gpu.yaml"
  user_account { username = "pesquisador"  keys = [ ... ] }
  ip_config { ipv4 { address = "dhcp" } }
}
\end{verbatim}
\end{minipage}
\medskip

O usuário e a chave de acesso são configurados pela API do Proxmox VE, enquanto os pacotes do sistema vêm do \textit{vendor-data}. Como o arquivo já está no \textit{storage}, o Terraform apenas o referencia: não há envio de arquivos nem acesso SSH ao host.

\subsection{Conteúdo do \textit{vendor-data} de GPU}

No primeiro \textit{boot}, o \textit{cloud-init} executa as ações abaixo (resumo do arquivo de \textit{vendor-data} de GPU):

\begin{itemize}
    \item instala e habilita o \textit{QEMU Guest Agent};
    \item instala o Docker Engine, o \textit{Buildx} e o \textit{Compose} a partir do repositório oficial, e adiciona os usuários da VM ao grupo \texttt{docker};
    \item instala o cliente NFS e monta em \texttt{/mnt/research} o diretório compartilhado de dados de pesquisa do laboratório;
    \item instala o \textit{driver} NVIDIA de módulo de \textit{kernel} aberto (\texttt{nvidia-driver-595-open});
    \item reinicia a VM ao final, para que o módulo do \textit{driver} seja carregado.
\end{itemize}

\medskip
\noindent\begin{minipage}{\textwidth}
\footnotesize
\begin{verbatim}
#cloud-config
package_update: true
packages: [qemu-guest-agent, ca-certificates, curl, gnupg,
           nfs-common, ubuntu-drivers-common]
runcmd:
  - systemctl enable --now qemu-guest-agent
  - ... instala docker-ce, docker-compose-plugin ...
  - ... monta NFS em /mnt/research ...
  - apt-get install -y nvidia-driver-595-open
power_state:
  mode: reboot
  condition: true
\end{verbatim}
\end{minipage}
\medskip

\subsection{Verificação e cuidados}

Como a VM com GPU nasce desligada, o \textit{cloud-init} executa quando o usuário a liga pela primeira vez. A instalação exige acesso à internet e leva alguns minutos, seguida de uma reinicialização; durante esse período a VM aparece como ligada, mas ainda não está pronta. Concluída a etapa, a instalação pode ser verificada na VM com o comando \texttt{nvidia-smi}. O \textit{snippet} deve ser mantido pelo administrador, pois a versão do \textit{driver} é fixada no arquivo. Também é necessário que o \textit{template} seja do tipo \texttt{q35} (seção anterior) e que a VM alcance o servidor NFS do laboratório.

	\end{appendix}
	
\end{document} 
